\documentclass[10pt,a4paper]{amsart}
\usepackage[margin=19mm]{geometry}
\usepackage{amsmath,amssymb,mathtools}
\usepackage[scr=rsfs,cal=euler]{mathalfa}
\usepackage{xfrac}
\usepackage[unicode,hidelinks,bookmarks=false]{hyperref}
\newcommand{\ie}{i.e.\ }
\newcommand{\cf}{cf.\ }
\newcommand{\resp}{respectively\ }
\newcommand{\Subsection}{\subsection}
\providecommand{\nobreakdash}{\nobreak\hbox{-}}
\setlength{\emergencystretch}{2em}
\multlinegap0pt
\DeclareMathOperator{\projective}{\mathbb{P}}
\usepackage{amssymb}
\newtheorem{proposition}{Proposition}
\newtheorem{lemma}{Lemma}
\DeclareMathOperator{\Ann}{Ann}
\DeclareMathOperator{\Ass}{Ass}
\DeclareMathOperator{\Ext}{Ext}
\DeclareMathOperator{\Supp}{Supp}
\DeclareMathOperator{\codim}{codim}
\DeclareMathOperator{\coker}{coker}
\DeclareMathOperator{\depth}{depth}
\newcommand{\frakp}{{\mathfrak p}}
\DeclareMathOperator{\height}{ht}
\DeclareMathOperator{\image}{Im}
\title{On codimension-two subcanonical varieties inside $\projective^n$}
\author{Manoj Kummini, Abhiram Subramanian}
\date{Source: arXiv:2511.19962v3 (CC BY 4.0)}
\begin{document}
\maketitle

\noindent\emph{Source and licence.} On codimension-two subcanonical varieties inside $\projective^n$. Version arXiv:2511.19962v3 (CC BY 4.0), distributed by arXiv under the Creative Commons Attribution 4.0 International licence, http://creativecommons.org/licenses/by/4.0/, as recorded in the arXiv licence metadata for this exact version and shown on the abstract page https://arxiv.org/abs/2511.19962v3. Full licence notice: \texttt{licenses/arxiv-2511.19962-CC-BY-4.0.txt}.\par\smallskip

\noindent\textit{Editorial scope.} Counted unit: the article's lemma lemma:Mmodxi (source0.tex lines 760--783). The article's complete original proof of it is delivered with it (source0.tex lines 785--842, one \texttt{proof} environment of 58 lines). No mathematics is written, repaired or re-derived: the statement and the proof are the article's own lines, in the article's own notation, and no retained line is substituted or re-flowed.  Retained uncounted as the source-local context the proof needs: lines 591--597; lines 599--604; lines 737--740.  Nothing was omitted from any retained span, so the package carries no omission record.  No retained line contains a citation, so the wrapper carries no bibliography and no bibliography entry.  No retained line contains a figure environment, an image-inclusion command or drawing code, so no image is delivered and the package declares no asset dependency.  Editorial formatting only, in addition: an AMS \texttt{amsart} preamble that restates the article's own declarations and macros used by the retained lines, namely the environments \texttt{proposition}, \texttt{lemma} on separate counters, together with the standard packages those lines need.  The retained environments print in this document as Proposition 1, Lemma 1 in order of appearance, whereas the article prints the same items with the numbering of its own full document.  The complete native source of the article as distributed in the arXiv e-print for this exact version is bundled unchanged as \texttt{sources/189750/source0.tex}.
\begin{proposition}
\label{proposition:SkAndExt}
Let $S$ be a regular local ring and $M$ a finitely generated $S$-module.
Let $k \geq 1$.
Then $M$ satisfies $(S_k)$ if and only if for all $i \geq 1$, $\codim_S
\Ext_S^i(M,S ) \geq i+k$.
\end{proposition}

\begin{lemma}
\label{lemma:S2UFDfree}
Let $S$ be a noetherian UFD and $M$ a finitely generated rank-one
$S$-module satisfying $(S_2)$. 
Then $M$ is free.
\end{lemma}

\begin{lemma}
\label{lemma:Mdual}
$M^* \simeq M(a_X +n + 1)$.
\end{lemma}

\begin{lemma}
\label{lemma:Mmodxi}
With notation as above, assume that $2d_1 \leq a_X +n + 2$. 
Further, write $i$ for the inclusion map $R\xi \to M$.
Let $\displaystyle J = \Ann_R\left(\coker \left(M^*
\stackrel{i^*}{\to}R(d_1)\right)\right)$.
Then 
\begin{enumerate}

\item
\label{lemma:Mmodxi:S1}
$M/R\xi$ satisfies $(S_1)$.

\item
\label{lemma:Mmodxi:NdblDual}
$(M/R\xi)^{**} = R(d_1-(a_X + n + 1))$.

\item
\label{lemma:Mmodxi:Junm}
$J$ is a unmixed height-two ideal and $J_\frakp$ is a complete intersection
for all $\frakp \neq R_+$.

\end{enumerate}
\end{lemma}

\begin{proof}
%
Write $N=M/R\xi$. Then we have an exact sequence
\[
0 \to N^* \to M^* \stackrel{i^*}\to R(d_1) \to \Ext^1_R(N,R) \to 
\Ext^1_R(M,R) \to 0
\]
and isomorphisms $\Ext^i_R(N,R) \stackrel\sim\to \Ext^i_R(M,R)$ for all $i
\geq 2$.
By Proposition~\ref{proposition:SkAndExt}, we need to show that 
$\codim_R(\Ext_R^i(N,R)) \geq i+1$.
Since $M$ satisfies $(S_2)$ and 
$\Ext_R^i(M,R )$ is a finite-length $R$-module for every $i>0$, 
it remains to show that $\codim_R(\Ext_R^1(N,R)) \geq 2$.
Since $\displaystyle\Supp \left(\Ext_R^1(N,R) \right ) = \Supp (R/J)$, we
need to show that $\height J \geq 2$.
It is immediate that $\height J >0$.
Note also that $N^*$ is free, since it is a rank-one module satisfying
$(S_2)$ (Lemma \ref{lemma:S2UFDfree}).

\eqref{lemma:Mmodxi:S1}. First assume that $J$ contains a linear form $l$.
If $\height J = 1$, then $J = (l)$, from which it would follow that
$\image i^*$ is a free $R$-module and that $M^*$ and $M$ are free.
This is a contradiction, so $\height J \geq 2$, which
proves~\eqref{lemma:Mmodxi:S1}.
Hence we may assume that $J$ does not contain a linear form, or,
equivalently, that $\image i^*$ is zero in degree $\leq -d_1+1$.

Since $M^* \simeq M(a_X +n + 1)$ (Lemma~\ref{lemma:Mdual}),
it is generated in degrees
$d_i - (a_X + n + 1)$, $1 \leq i \leq m$ and $0$.
Since $2d_1 \leq a_X + n + 2$, we see that
$d_1 - (a_X + n + 1) \leq -d_1+1$.
Therefore $N^*$ maps on to a minimal generator $\eta$ of $M^*$ degree 
$d_1 - (a_X + n + 1)$ and the remaining minimal generators of $M^*$ 
have degree $\geq -d_1+2$. 
It follows that 
$M^*/R\eta = \image i^* \simeq N(a_X+n+1)$, so $N$ satisfies $(S_1)$.

\eqref{lemma:Mmodxi:NdblDual}:
$N^* = R(a_X+n+1-d_1)$, by the above argument. 

\eqref{lemma:Mmodxi:Junm}:
By~\eqref{lemma:Mmodxi:S1}, $\height J \geq 2$.
Let $\frakp \in \Ass R/J$. 
Since $\depth M^* \geq 2$ and $\depth R \geq 3$, 
we see that $\frakp \neq R_+$.
Hence we have a free resolution
\[
0 \to N_\frakp^* \to M_\frakp^* \to R_\frakp \to  0
\]
of $(R/J)_{\frakp}$ over $R_\frakp$. Since $\depth_{R_\frakp } 
(R/J)_{\frakp} = 0$, we see from the Auslander-Buchsbaum formula that
$\height \frakp \leq 2$. Since $\height J \geq 2$, we conclude that
$\height \frakp = 2$. Hence $J$ is a unmixed height-two ideal.
Further for all $\frakp \neq R_+$, $M_\frakp^*$ is a rank-two free
$R_\frakp$-module, so $J_\frakp$ is a complete intersection.
\end{proof}

\end{document}
